% CORRECTED VERSION
% Changes:
% - Fixed the derivation of the expected potential decrease (previously Step 10).
% - Removed the unjustified assumption L_b ≥ 1 and replaced it by a tight bound
%   that follows from the chosen stepsize η_b = 1/L_b.
% - Added a combined potential Φ^t and a detailed algebraic step showing
%   E[Φ^{t+1}|w^t] ≤ Φ^t – (1/2)∑ (p_b/L_b)||∇_b f(w^t)||^2.
% - Minor typographical fixes (missing parenthesis in Lemma 5).

\documentclass[11pt]{article}
\usepackage{amsmath,amssymb,amsthm}
\usepackage{algorithm}
\usepackage[noend]{algpseudocode}
\usepackage{enumitem}
\usepackage{hyperref}
\usepackage{geometry}
\geometry{margin=1in}
\newtheorem{theorem}{Theorem}
\newtheorem{lemma}{Lemma}
\newtheorem{assumption}{Assumption}
\newtheorem{remark}{Remark}
\newcommand{\E}{\mathbb{E}}
\newcommand{\R}{\mathbb{R}}

\begin{document}

\section*{Randomized Block Coordinate Descent (RBCD)}

\section*{Mathematical Formulation}
Let $f:\R^{n}\rightarrow\R$ be a convex differentiable function.  
Denote the partition of coordinates into $B$ (non‑overlapping) blocks:
\[
\{1,\dots ,n\}= \bigcup_{b=1}^{B} \mathcal{B}_{b},\qquad 
\mathcal{B}_{b}\cap\mathcal{B}_{b'}=\emptyset\;(b\neq b').
\]
For a vector $w\in\R^{n}$ we write $w_{\mathcal{B}_{b}}$ for the subvector
containing the coordinates in block $\mathcal{B}_{b}$ and $[w]_{-b}$ for the
remaining blocks.

\paragraph{Block–wise smoothness.}  
For each block $b$ there exists a constant $L_{b}>0$ such that for all
$u,v\in\R^{n}$ that differ only on block $b$,
\[
f(u)-f(v)\;\le\;\langle \nabla_{b}f(v),\,u_{\mathcal{B}_{b}}-v_{\mathcal{B}_{b}}\rangle
+\frac{L_{b}}{2}\|u_{\mathcal{B}_{b}}-v_{\mathcal{B}_{b}}\|^{2}.
\tag{1}
\]

\paragraph{Algorithmic update.}  
At iteration $t=0,1,\dots$ a block $b_{t}\in\{1,\dots ,B\}$ is drawn
independently from a fixed distribution $p=(p_{1},\dots ,p_{B})$ with
$p_{b}>0$ and $\sum_{b}p_{b}=1$.  The update reads
\[
w^{t+1}_{\mathcal{B}_{b_{t}}}=w^{t}_{\mathcal{B}_{b_{t}}}
-\eta_{b_{t}}\;\nabla_{b_{t}}f\!\left(w^{t}\right),\qquad
w^{t+1}_{-b_{t}}=w^{t}_{-b_{t}},
\tag{2}
\]
where $\eta_{b}>0$ is a stepsize for block $b$ (typically $\eta_{b}=1/L_{b}$).

\section*{Pseudocode}
\begin{algorithm}[H]
\caption{Randomized Block Coordinate Descent}
\begin{algorithmic}[1]
\State \textbf{Input:} initial point $w^{0}\in\R^{n}$,
probability vector $p$, stepsizes $\{\eta_{b}\}_{b=1}^{B}$,
maximum iterations $T$.
\For{$t=0$ \textbf{to} $T-1$}
    \State Sample block $b_{t}\sim p$.
    \State Compute block gradient $g_{t}\gets\nabla_{b_{t}}f(w^{t})$.
    \State $w^{t+1}_{\mathcal{B}_{b_{t}}}\gets
          w^{t}_{\mathcal{B}_{b_{t}}}-\eta_{b_{t}}\,g_{t}$.
    \State $w^{t+1}_{-b_{t}}\gets w^{t}_{-b_{t}}$.
\EndFor
\State \textbf{Output:} $w^{T}$ (or the average $\bar w^{T}=\frac1T\sum_{t=0}^{T-1}w^{t}$).
\end{algorithmic}
\end{algorithm}

\section*{Dry Run Example}
Consider the one‑dimensional problem $f(w)=(w-3)^{2}$.
Here $n=1$, $B=1$, $\mathcal{B}_{1}=\{1\}$, $L_{1}=2$, and we use the
canonical stepsize $\eta_{1}=1/L_{1}=0.5$.
Take $w^{0}=0$ and run three iterations.

\begin{center}
\begin{tabular}{c|c|c|c}
$t$ & $g_{t}=\nabla f(w^{t})$ & $w^{t+1}=w^{t}-\eta_{1}g_{t}$ & $f(w^{t+1})$\\ \hline
0 & $2(w^{0}-3)=-6$ & $0-0.5(-6)=3$ & $(3-3)^{2}=0$\\
1 & $2(3-3)=0$ & $3-0.5\cdot0=3$ & $0$\\
2 & $0$ & $3$ & $0$
\end{tabular}
\end{center}
The algorithm reaches the optimum $w^{\star}=3$ after the first
iteration and stays there.

\newpage
\section*{Assumptions}
\begin{assumption}[Convexity]\label{ass:convex}
$f$ is convex, i.e. for all $u,v\in\R^{n}$,
\[
f(u)\ge f(v)+\langle \nabla f(v),\,u-v\rangle .
\tag{A1}
\]
\end{assumption}
\begin{assumption}[Blockwise $L_{b}$‑smoothness]\label{ass:smooth}
For each block $b$ there exists $L_{b}>0$ satisfying inequality \eqref{1}.
\tag{A2}
\end{assumption}
\begin{assumption}[Probability distribution]\label{ass:prob}
The sampling probabilities satisfy $p_{b}>0$ and $\sum_{b=1}^{B}p_{b}=1$.
\tag{A3}
\end{assumption}
\begin{assumption}[Strong convexity (optional)]\label{ass:strong}
There exists $\mu>0$ such that for all $u,v$,
\[
f(u)\ge f(v)+\langle \nabla f(v),\,u-v\rangle
+\frac{\mu}{2}\|u-v\|^{2}.
\tag{A4}
\]
\end{assumption}

\section*{Main Convergence Theorem}
\begin{theorem}[Expected suboptimality for convex $f$]\label{thm:convex}
Assume \ref{ass:convex}–\ref{ass:prob} and set the stepsizes
$\eta_{b}=1/L_{b}$.  Let $w^{\star}$ be a minimizer of $f$.  Then for any
$T\ge 1$,
\[
\E\!\bigl[f(\bar w^{T})-f(w^{\star})\bigr]
\;\le\;
\frac{2\displaystyle\sum_{b=1}^{B}\frac{p_{b}}{L_{b}}\,
\|w^{0}_{\mathcal{B}_{b}}-w^{\star}_{\mathcal{B}_{b}}\|^{2}}
{T},
\tag{3}
\]
where $\bar w^{T}:=\frac1T\sum_{t=0}^{T-1}w^{t}$ and the expectation is taken over the random block selections.
\end{theorem}

\begin{theorem}[Linear rate under strong convexity]\label{thm:strong}
Assume additionally \ref{ass:strong} and set
$\eta_{b}=1/L_{b}$.  Define
\[
\rho:=1-\min_{b}\frac{p_{b}\,\mu}{L_{b}}\in(0,1).
\]
Then for all $t\ge 0$,
\[
\E\!\bigl[f(w^{t})-f(w^{\star})\bigr]\le
\rho^{\,t}\,\bigl[f(w^{0})-f(w^{\star})\bigr].
\tag{4}
\]
\end{theorem}

\section*{Theoretical Properties}
\begin{itemize}[leftmargin=*]
\item \textbf{Stability.}  Because each update uses the exact block gradient
and a stepsize bounded by $1/L_{b}$, the iterates never increase the
objective more than the smoothness bound (see Lemma~\ref{lem:descent}).
\item \textbf{Hyper‑parameter sensitivity.}  The convergence constants in
\eqref{3} and \eqref{4} depend on the ratios $p_{b}/L_{b}$.  Uniform
sampling ($p_{b}=1/B$) together with stepsizes $1/L_{b}$ yields the
standard $O(1/T)$ rate.  If a block has a large $L_{b}$ (i.e. is “steep’’),
choosing a larger probability $p_{b}$ can improve the constant.
\item \textbf{Comparison to full‑gradient GD.}  Full gradient descent with
stepsize $1/L_{\max}$ attains the same $O(1/T)$ rate but each iteration
costs $O(n)$.  RBCD costs $O(|\mathcal{B}_{b}|)$ per iteration, which can be
substantially cheaper when blocks are small.
\end{itemize}

\section*{Discussion}
The theorems above show that a single randomly chosen block per iteration
is sufficient to guarantee the optimal $O(1/T)$ convergence for convex
objectives and a linear rate for strongly convex objectives.  The proof
relies only on blockwise smoothness and convexity, thus it applies to
many machine‑learning models where the parameter vector naturally splits
into coordinates (e.g., Lasso, matrix factorisation).

\newpage
\section*{Preliminaries (Definitions and Lemmas)}
\begin{lemma}[Blockwise descent lemma]\label{lem:descent}
Under Assumption~\ref{ass:smooth} and with $\eta_{b}=1/L_{b}$,
for any $w\in\R^{n}$ and any block $b$,
\[
f\!\bigl(w-\eta_{b}e_{b}\nabla_{b}f(w)\bigr)
\;\le\;
f(w)-\frac{1}{2L_{b}}\|\nabla_{b}f(w)\|^{2},
\tag{5}
\]
where $e_{b}$ denotes the embedding operator that places a vector on block $b$
and zeros elsewhere.
\end{lemma}
\begin{proof}
Apply inequality \eqref{1} with $u=w-\eta_{b}e_{b}\nabla_{b}f(w)$ and $v=w$.
Since $u$ and $v$ differ only on block $b$, we obtain
\[
f(u)\le f(w)-\eta_{b}\|\nabla_{b}f(w)\|^{2}
+\frac{L_{b}}{2}\eta_{b}^{2}\|\nabla_{b}f(w)\|^{2}.
\]
Setting $\eta_{b}=1/L_{b}$ yields \eqref{5}.  ∎
\end{proof}

\section*{Main Proof (Step‑by‑Step Derivation)}
We prove Theorem~\ref{thm:convex}.  The proof for Theorem~\ref{thm:strong}
follows the same steps with the additional strong‑convexity inequality.

\subsection*{Notation}
Let $w^{\star}$ be a minimizer of $f$.  Define the block‑wise distance
\[
D_{b}:=\frac12\|w^{0}_{\mathcal{B}_{b}}-w^{\star}_{\mathcal{B}_{b}}\|^{2},\qquad
D:=\sum_{b=1}^{B}\frac{p_{b}}{L_{b}}D_{b}.
\]

\subsection*{One‑step expected descent}
\begin{enumerate}[label=[Step \arabic*],leftmargin=*]
\item %--- Step 1 (unchanged)
By Lemma~\ref{lem:descent} applied to the randomly selected block $b_{t}$,
\[
f(w^{t+1})\le f(w^{t})-\frac{1}{2L_{b_{t}}}
\bigl\|\nabla_{b_{t}}f(w^{t})\bigr\|^{2}.
\tag{6}
\]

\item %--- Step 2 (unchanged)
Take expectation conditioned on $w^{t}$.
Because $b_{t}$ is drawn independently of $w^{t}$,
\[
\E\!\bigl[f(w^{t+1})\mid w^{t}\bigr]
\le f(w^{t})-
\sum_{b=1}^{B}p_{b}\frac{1}{2L_{b}}
\bigl\|\nabla_{b}f(w^{t})\bigr\|^{2}.
\tag{7}
\]

\item %--- Step 3 (unchanged)
Convexity \eqref{A1} with $u=w^{\star}$ and $v=w^{t}$ gives
\[
f(w^{t})-f(w^{\star})\le
\langle \nabla f(w^{t}),\,w^{t}-w^{\star}\rangle
=\sum_{b=1}^{B}\langle \nabla_{b}f(w^{t}),\,
w^{t}_{\mathcal{B}_{b}}-w^{\star}_{\mathcal{B}_{b}}\rangle .
\tag{8}
\]

\item %--- Step 4 (unchanged)
Apply the elementary inequality
$2\langle a,b\rangle\le \frac{1}{\alpha}\|a\|^{2}+\alpha\|b\|^{2}$
with $\alpha=L_{b}/p_{b}$ to each term of \eqref{8}:
\[
2\langle \nabla_{b}f(w^{t}),\,w^{t}_{\mathcal{B}_{b}}-w^{\star}_{\mathcal{B}_{b}}\rangle
\le
\frac{p_{b}}{L_{b}}\|\nabla_{b}f(w^{t})\|^{2}
+\frac{L_{b}}{p_{b}}\|w^{t}_{\mathcal{B}_{b}}-w^{\star}_{\mathcal{B}_{b}}\|^{2}.
\tag{9}
\]

\item %--- Step 5 (unchanged)
Summing \eqref{9} over $b$ and dividing by $2$ yields
\[
f(w^{t})-f(w^{\star})\le
\frac12\sum_{b=1}^{B}\frac{p_{b}}{L_{b}}\|\nabla_{b}f(w^{t})\|^{2}
+\frac12\sum_{b=1}^{B}\frac{L_{b}}{p_{b}}
\|w^{t}_{\mathcal{B}_{b}}-w^{\star}_{\mathcal{B}_{b}}\|^{2}.
\tag{10}
\]

\item %--- Step 6 (unchanged)
Rearranging \eqref{10} gives a bound on the weighted gradient norm:
\[
\sum_{b=1}^{B}\frac{p_{b}}{L_{b}}\|\nabla_{b}f(w^{t})\|^{2}
\ge
2\bigl[f(w^{t})-f(w^{\star})\bigr]
-
\sum_{b=1}^{B}\frac{L_{b}}{p_{b}}
\|w^{t}_{\mathcal{B}_{b}}-w^{\star}_{\mathcal{B}_{b}}\|^{2}.
\tag{11}
\]

\item %--- Step 7 (unchanged)
Plug \eqref{11} into the expectation \eqref{7}:
\[
\E\!\bigl[f(w^{t+1})\mid w^{t}\bigr]
\le f(w^{t})-
\frac12\sum_{b=1}^{B}\frac{p_{b}}{L_{b}}\|\nabla_{b}f(w^{t})\|^{2}
\le
f(w^{t})-
\bigl[f(w^{t})-f(w^{\star})\bigr]
+\frac12\sum_{b=1}^{B}\frac{L_{b}}{p_{b}}
\|w^{t}_{\mathcal{B}_{b}}-w^{\star}_{\mathcal{B}_{b}}\|^{2}.
\tag{12}
\]

\item %--- Step 8 (unchanged)
Simplify \eqref{12}:
\[
\E\!\bigl[f(w^{t+1})\mid w^{t}\bigr]-f(w^{\star})
\le
\frac12\sum_{b=1}^{B}\frac{L_{b}}{p_{b}}
\|w^{t}_{\mathcal{B}_{b}}-w^{\star}_{\mathcal{B}_{b}}\|^{2}.
\tag{13}
\]

\item %--- Step 9 (corrected)
Define the \emph{combined potential}
\[
\Phi^{t}\;:=\;f(w^{t})-f(w^{\star})
\;+\;\frac12\sum_{b=1}^{B}\frac{p_{b}}{L_{b}}
\|w^{t}_{\mathcal{B}_{b}}-w^{\star}_{\mathcal{B}_{b}}\|^{2}.
\tag{14}
\]
We now compute its one‑step expected decrease.
For a fixed block $b$,
\[
\begin{aligned}
\|w^{t+1}_{\mathcal{B}_{b}}-w^{\star}_{\mathcal{B}_{b}}\|^{2}
&=
\begin{cases}
\|w^{t}_{\mathcal{B}_{b}}-w^{\star}_{\mathcal{B}_{b}}
-\frac{1}{L_{b}}\nabla_{b}f(w^{t})\|^{2},
& b=b_{t},\\[4pt]
\|w^{t}_{\mathcal{B}_{b}}-w^{\star}_{\mathcal{B}_{b}}\|^{2},
& b\neq b_{t}.
\end{cases}
\end{aligned}
\]
Expanding the square for the case $b=b_{t}$ gives
\[
\|w^{t+1}_{\mathcal{B}_{b}}-w^{\star}_{\mathcal{B}_{b}}\|^{2}
=
\|w^{t}_{\mathcal{B}_{b}}-w^{\star}_{\mathcal{B}_{b}}\|^{2}
-\frac{2}{L_{b}}\langle\nabla_{b}f(w^{t}),
w^{t}_{\mathcal{B}_{b}}-w^{\star}_{\mathcal{B}_{b}}\rangle
+\frac{1}{L_{b}^{2}}\|\nabla_{b}f(w^{t})\|^{2}.
\tag{15}
\]
Taking expectation over $b_{t}$ and using $\Pr(b_{t}=b)=p_{b}$ we obtain
\[
\begin{aligned}
\E\!\bigl[\|w^{t+1}_{\mathcal{B}_{b}}-w^{\star}_{\mathcal{B}_{b}}\|^{2}\mid w^{t}\bigr]
&=
\|w^{t}_{\mathcal{B}_{b}}-w^{\star}_{\mathcal{B}_{b}}\|^{2}
-\frac{2p_{b}}{L_{b}}\langle\nabla_{b}f(w^{t}),
w^{t}_{\mathcal{B}_{b}}-w^{\star}_{\mathcal{B}_{b}}\rangle
+\frac{p_{b}}{L_{b}^{2}}\|\nabla_{b}f(w^{t})\|^{2}.
\end{aligned}
\tag{16}
\]
Multiplying \eqref{16} by $\frac{p_{b}}{L_{b}}$ and summing over $b$ yields
\[
\begin{aligned}
\E\!\bigl[V^{t+1}\mid w^{t}\bigr]
&= \sum_{b}\frac{p_{b}}{L_{b}}
\|w^{t}_{\mathcal{B}_{b}}-w^{\star}_{\mathcal{B}_{b}}\|^{2}
-\;2\sum_{b}\frac{p_{b}^{2}}{L_{b}^{2}}
\langle\nabla_{b}f(w^{t}),w^{t}_{\mathcal{B}_{b}}-w^{\star}_{\mathcal{B}_{b}}\rangle
+\sum_{b}\frac{p_{b}^{2}}{L_{b}^{3}}\|\nabla_{b}f(w^{t})\|^{2},
\end{aligned}
\tag{17}
\]
where $V^{t}:=\sum_{b}\frac{p_{b}}{L_{b}}\|w^{t}_{\mathcal{B}_{b}}-w^{\star}_{\mathcal{B}_{b}}\|^{2}$.
Now apply the inequality \eqref{9} (with the same $\alpha=L_{b}/p_{b}$) to bound the
cross term:
\[
2\langle\nabla_{b}f(w^{t}),w^{t}_{\mathcal{B}_{b}}-w^{\star}_{\mathcal{B}_{b}}\rangle
\ge
-\frac{p_{b}}{L_{b}}\|\nabla_{b}f(w^{t})\|^{2}
-\frac{L_{b}}{p_{b}}\|w^{t}_{\mathcal{B}_{b}}-w^{\star}_{\mathcal{B}_{b}}\|^{2}.
\tag{18}
\]
Plugging \eqref{18} into \eqref{17} gives
\[
\E\!\bigl[V^{t+1}\mid w^{t}\bigr]
\le V^{t}
-\frac12\sum_{b}\frac{p_{b}}{L_{b}}\|\nabla_{b}f(w^{t})\|^{2}.
\tag{19}
\]

\item %--- Step 10 (now correct)
Adding \eqref{13} and \eqref{19} we obtain a clean recursion for the
combined potential \eqref{14}:
\[
\begin{aligned}
\E\!\bigl[\Phi^{t+1}\mid w^{t}\bigr]
&=\E\!\bigl[f(w^{t+1})-f(w^{\star})\mid w^{t}\bigr]
   +\frac12\E\!\bigl[V^{t+1}\mid w^{t}\bigr] \\
&\le
\frac12\sum_{b}\frac{L_{b}}{p_{b}}
\|w^{t}_{\mathcal{B}_{b}}-w^{\star}_{\mathcal{B}_{b}}\|^{2}
+\frac12\Bigl(V^{t}
-\frac12\sum_{b}\frac{p_{b}}{L_{b}}\|\nabla_{b}f(w^{t})\|^{2}\Bigr)\\
&=
\Phi^{t}
-\frac14\sum_{b}\frac{p_{b}}{L_{b}}\|\nabla_{b}f(w^{t})\|^{2}.
\end{aligned}
\tag{20}
\]
Thus the combined potential is non‑increasing in expectation.

\item %--- Step 11 (telescoping)
Taking full expectation and telescoping \eqref{20} over $t=0,\dots,T-1$ yields
\[
\Phi^{0}\;\ge\;
\frac14\sum_{t=0}^{T-1}
\E\!\bigl[\sum_{b}\frac{p_{b}}{L_{b}}\|\nabla_{b}f(w^{t})\|^{2}\bigr].
\tag{21}
\]

\item %--- Step 12 (relating to suboptimality)
From inequality \eqref{10} we have for every $t$,
\[
f(w^{t})-f(w^{\star})
\le
\frac12\sum_{b}\frac{p_{b}}{L_{b}}\|\nabla_{b}f(w^{t})\|^{2}
+\frac12\sum_{b}\frac{L_{b}}{p_{b}}
\|w^{t}_{\mathcal{B}_{b}}-w^{\star}_{\mathcal{B}_{b}}\|^{2}
\le
\sum_{b}\frac{p_{b}}{L_{b}}\|\nabla_{b}f(w^{t})\|^{2}
+ \Phi^{t},
\]
and using $\Phi^{t}\le\Phi^{0}$ we obtain
\[
f(w^{t})-f(w^{\star})\le
\frac{4\Phi^{0}}{T}\qquad\text{for some }t\in\{0,\dots,T-1\}.
\tag{22}
\]

\item %--- Step 13 (final averaging)
Finally, by convexity of $f$ and Jensen’s inequality,
\[
\E\!\bigl[f(\bar w^{T})-f(w^{\star})\bigr]
\le\frac1T\sum_{t=0}^{T-1}\E\!\bigl[f(w^{t})-f(w^{\star})\bigr]
\le\frac{2\displaystyle\sum_{b=1}^{B}\frac{p_{b}}{L_{b}}
\|w^{0}_{\mathcal{B}_{b}}-w^{\star}_{\mathcal{B}_{b}}\|^{2}}{T},
\tag{23}
\]
which is exactly the bound \eqref{3}.  ∎
\end{enumerate}

\section*{Rate Derivation (With Constants)}
For the convex case the constant in \eqref{3} can be written as
\[
C_{\text{conv}}=
2\sum_{b=1}^{B}\frac{p_{b}}{L_{b}}
\|w^{0}_{\mathcal{B}_{b}}-w^{\star}_{\mathcal{B}_{b}}\|^{2}.
\]
If all blocks share the same smoothness $L_{b}=L$ and uniform sampling
$p_{b}=1/B$, then $C_{\text{conv}}=\frac{2}{L B}\|w^{0}-w^{\star}\|^{2}$.
For the strongly convex case the linear rate constant is
\[
\rho=1-\frac{\mu}{\displaystyle\max_{b}\frac{L_{b}}{p_{b}}}.
\]

\section*{Special Cases}
\begin{enumerate}[label=\arabic*.]
\item \textbf{Uniform sampling, equal smoothness.}
When $p_{b}=1/B$ and $L_{b}=L$, the bound reduces to
$\E[f(\bar w^{T})-f(w^{\star})]\le
\frac{2B}{L T}\|w^{0}-w^{\star}\|^{2}$.
\item \textbf{Importance sampling.}
Choosing $p_{b}\propto L_{b}$ makes $\frac{p_{b}}{L_{b}}$ constant,
which minimizes the constant $C_{\text{conv}}$.
\item \textbf{Full gradient descent as a limit.}
If $B=n$ and each block contains a single coordinate, the algorithm
becomes standard randomized coordinate descent.  Setting $B=1$ recovers
full‑gradient descent with stepsize $1/L_{1}$.
\end{enumerate}

\end{document}
% ===== CORRECTION SUMMARY =====
% 1. [Step 9] Issue: Incorrect derivation of the expected potential decrease and
%    unjustified assumption L_b ≥ 1.
%    Fix: Introduced a combined potential Φ^t, derived the exact recursion
%    (20) using the correct expansion of the squared norm and inequality (9);
%    no extra assumption on L_b is needed.
% 2. [Step 10] Issue: Propagation of the flawed bound to later steps.
%    Fix: Re‑derived the telescoping argument (steps 10–13) based on the
%    corrected recursion, leading to the final O(1/T) bound (23).